% INDEPENDENT INSTRUCTIONAL MATERIAL --- NOT WEN-WEI LI'S SOURCE TEXT.
% stable-unit-id: o014.mastery.diagram-chasing.001
% provenance: OpenAI Codex gpt-5.6-sol, Ultra
% license: CC BY 4.0; compatible with the source translation license.

\section*{Mastery Bridge 001: Diagram Chasing}
\addcontentsline{toc}{section}{Mastery Bridge 001: Diagram Chasing}
\hypertarget{o014.mastery.diagram-chasing.001}{ }
\label{mastery:diagram-chasing-001}

\noindent\textbf{Status and attribution.}
This unit is independent instructional material prepared to accompany this
book. It is not part of Wen-Wei Li's source text and must not be attributed to
him. Production provenance:
\texttt{OpenAI Codex gpt-5.6-sol, Ultra}. This material is available under the
Creative Commons Attribution 4.0 International license
(CC BY 4.0; \url{https://creativecommons.org/licenses/by/4.0/}), and is
therefore compatible with the license of the source text.

\subsection*{Prerequisite map}

This unit can be read after mastering the following short chain of
prerequisites.
\begin{compactitem}
	\item Kernels, cokernels, and their functoriality:
	\S\ref{sec:Ker-Coker}, especially
	\eqref{eqn:Ker-Coker-functoriality}.
	\item Images, coimages, and epi--mono factorization:
	\S\ref{sec:images}.
	\item Abelian categories and the stability of epimorphisms and
	monomorphisms under pullback and pushout:
	\S\ref{sec:Abel-cat-def} and
	Proposition~\ref{prop:Abel-cat-pull-push}.
	\item Complexes, cohomology, and exactness:
	\S\ref{sec:cohomology}, especially
	Definition~\ref{def:cohomology}.
	\item The categorical construction of connecting morphisms and the Snake
	Lemma:
	\S\ref{sec:diagram-lemmas}, Remark~\ref{rem:connecting-canonical}, and
	Theorem~\ref{prop:snake-lemma}.
\end{compactitem}
Conceptually, the progression is
\[
	\text{kernels/cokernels}
	\Longrightarrow \text{images and exactness}
	\Longrightarrow \text{lift--map--descend}
	\Longrightarrow \text{connecting morphisms}.
\]

\subsection*{Working method: start with the obstruction}

For element chasing in a module category, use the following discipline.
\begin{enumerate}
	\item State precisely what must be proved: membership in a kernel or
	image, or equality of two cokernel classes.
	\item Start with the element that measures the obstruction, usually a
	kernel element or a cokernel class.
	\item Lift an element only along an arrow known to be surjective. Map an
	element forward along any available arrow.
	\item Use commutativity to move an equation to the other side of the
	diagram, then use exactness to replace a statement of ``being zero'' by
	``coming from the preceding image.''
	\item If a choice is required, prove that the difference between any two
	choices vanishes in the final target.
\end{enumerate}
In a general abelian category, an ``element'' can be replaced by a morphism
from a test object. When a lift is not directly available, use an epimorphic
cover as in Lemma~\ref{prop:Abel-cat-exact-aux}. In this way, the element
chases below are not merely heuristic: they are concrete models for arguments
with kernels, cokernels, and factorizations.

As a worked example, consider the commutative diagram of $R$-modules
\[
\begin{tikzcd}
	0 \arrow[r] & A' \arrow[r, "i"] \arrow[d, "a'"'] & A \arrow[r, "p"] \arrow[d, "a"'] & A'' \arrow[r] \arrow[d, "a''"] & 0 \\
	0 \arrow[r] & B' \arrow[r, "j"'] & B \arrow[r, "q"'] & B'' \arrow[r] & 0
\end{tikzcd}
\]
with both rows exact. Suppose that $a'$ and $a''$ are injective. To prove
that $a$ is injective, start with the obstruction $x\in\Ker(a)$.
Commutativity gives $a''p(x)=qa(x)=0$, so the injectivity of $a''$ gives
$p(x)=0$. Exactness of the top row gives $x=i(x')$ for some $x'\in A'$.
Next,
\[
	0=a(x)=a i(x')=j a'(x').
\]
Because $j$ and $a'$ are injective, $x'=0$, and hence $x=0$. Notice the
recurring pattern: start in a kernel, move right, use injectivity, return to
the left by exactness, and then use injectivity once more.

\subsection*{Exercises and complete solutions}

% stable-exercise-id: o014.mastery.diagram-chasing.001.ex01
\subsubsection*{Exercise 1 --- Functoriality of kernels, cokernels, and images}
\hypertarget{o014.mastery.diagram-chasing.001.ex01}{ }
\label{mastery:diagram-chasing-001-ex01}

Let
\[
\begin{tikzcd}
	A \arrow[r, "f"] \arrow[d, "a"'] & B \arrow[d, "b"] \\
	A' \arrow[r, "g"'] & B'
\end{tikzcd}
\qquad (ga=bf)
\]
be a commutative square of $R$-modules.
\begin{enumerate}[(i)]
	\item Construct the natural homomorphisms
	$\overline a:\Ker(f)\to\Ker(g)$ and
	$\overline b:\Coker(f)\to\Coker(g)$.
	\item Prove that $b(\Image(f))\subset\Image(g)$, with equality if $a$
	is surjective.
	\item Prove that $\Ker(f)=a^{-1}(\Ker(g))$ if $b$ is injective.
\end{enumerate}

% stable-solution-id: o014.mastery.diagram-chasing.001.sol01
\paragraph{Solution.}
\hypertarget{o014.mastery.diagram-chasing.001.sol01}{ }
If $x\in\Ker(f)$, then $g(a(x))=b(f(x))=0$, so
$a(x)\in\Ker(g)$. Restricting $a$ gives
$\overline a:\Ker(f)\to\Ker(g)$. For cokernels, define
\[
	\overline b([y])=[b(y)], \qquad [y]\in B/\Image(f).
\]
If $y$ is replaced by $y+f(x)$, then
$b(y+f(x))=b(y)+g(a(x))$, which determines the same class modulo
$\Image(g)$. Thus $\overline b$ is well defined. These two constructions
are precisely those in \eqref{eqn:Ker-Coker-functoriality}.

Next,
\[
	b(\Image(f))=b(f(A))=g(a(A))\subset g(A')=\Image(g).
\]
If $a$ is surjective, then $a(A)=A'$, so the inclusion is an equality.
Finally, $x\in\Ker(f)$ always implies $a(x)\in\Ker(g)$. Conversely, if
$a(x)\in\Ker(g)$, then $0=g(a(x))=b(f(x))$; injectivity of $b$ gives
$f(x)=0$. Thus $x\in\Ker(f)$, and consequently
$\Ker(f)=a^{-1}(\Ker(g))$.

% stable-exercise-id: o014.mastery.diagram-chasing.001.ex02
\subsubsection*{Exercise 2 --- Intersections, sums, and the isomorphism theorem}
\hypertarget{o014.mastery.diagram-chasing.001.ex02}{ }
\label{mastery:diagram-chasing-001-ex02}

Let $U$ and $V$ be submodules of an $R$-module $M$. Define
\[
	r:U\cap V\longrightarrow U\oplus V,\quad r(w)=(w,-w),
	\qquad
	s:U\oplus V\longrightarrow U+V,\quad s(u,v)=u+v.
\]
Prove that
\[
	0\longrightarrow U\cap V\xrightarrow{r}U\oplus V
	\xrightarrow{s}U+V\longrightarrow0
\]
is exact. Use a kernel-and-image chase to obtain the canonical isomorphism
\[
	U/(U\cap V)\rightiso (U+V)/V.
\]

% stable-solution-id: o014.mastery.diagram-chasing.001.sol02
\paragraph{Solution.}
\hypertarget{o014.mastery.diagram-chasing.001.sol02}{ }
If $r(w)=0$, then the first component gives $w=0$, so $r$ is injective.
Moreover, $s(r(w))=w-w=0$, and hence $\Image(r)\subset\Ker(s)$. If
$(u,v)\in\Ker(s)$, then $u+v=0$. Thus $u=-v$ belongs to both $U$ and $V$,
and
\[
	(u,v)=(u,-u)=r(u).
\]
Consequently $\Ker(s)=\Image(r)$. Every element of $U+V$ has the form
$u+v$, so $s$ is surjective. The sequence is exact.

Now consider the homomorphism
\[
	\theta:U\longrightarrow (U+V)/V,\qquad u\longmapsto u+V.
\]
This map is surjective because $(u+v)+V=u+V$. Its kernel consists of all
$u\in U$ that also belong to $V$, namely $U\cap V$. Epi--mono
factorization, or the Isomorphism Theorem~\ref{prop:Abel-cat-isom-thm},
gives the canonical isomorphism
$U/(U\cap V)\rightiso (U+V)/V$.

% stable-exercise-id: o014.mastery.diagram-chasing.001.ex03
\subsubsection*{Exercise 3 --- Short Five Lemma}
\hypertarget{o014.mastery.diagram-chasing.001.ex03}{ }
\label{mastery:diagram-chasing-001-ex03}

In the commutative diagram
\[
\begin{tikzcd}
	0 \arrow[r] & A' \arrow[r, "i"] \arrow[d, "a'"'] & A \arrow[r, "p"] \arrow[d, "a"'] & A'' \arrow[r] \arrow[d, "a''"] & 0 \\
	0 \arrow[r] & B' \arrow[r, "j"'] & B \arrow[r, "q"'] & B'' \arrow[r] & 0,
\end{tikzcd}
\]
suppose that both rows are exact.
\begin{enumerate}[(i)]
	\item If $a'$ and $a''$ are injective, prove that $a$ is injective.
	\item If $a'$ and $a''$ are surjective, prove that $a$ is surjective.
	\item Conclude that if $a'$ and $a''$ are isomorphisms, then $a$ is
	also an isomorphism.
\end{enumerate}

% stable-solution-id: o014.mastery.diagram-chasing.001.sol03
\paragraph{Solution.}
\hypertarget{o014.mastery.diagram-chasing.001.sol03}{ }
For (i), take $x\in A$ with $a(x)=0$. Then
$a''p(x)=qa(x)=0$. Since $a''$ is injective, $p(x)=0$; exactness of the
top row gives $x=i(x')$. Next,
$0=a i(x')=j a'(x')$. Since $j$ and $a'$ are injective, $x'=0$, and hence
$x=0$.

For (ii), take $y\in B$ and set $y''=q(y)$. Surjectivity of $a''$ gives
$x''\in A''$ with $a''(x'')=y''$. Since $p$ is surjective, choose
$x\in A$ with $p(x)=x''$. Now
\[
	q(y-a(x))=y''-a''p(x)=0.
\]
Exactness of the bottom row gives $y-a(x)=j(y')$ for some $y'\in B'$.
Surjectivity of $a'$ gives $x'\in A'$ with $a'(x')=y'$. Therefore
\[
	y=a(x)+j(a'(x'))=a(x+i(x')).
\]
Thus $a$ is surjective. If the two outer morphisms are isomorphisms, (i)
and (ii) show that $a$ is both injective and surjective, and therefore an
isomorphism.

% stable-exercise-id: o014.mastery.diagram-chasing.001.ex04
\subsubsection*{Exercise 4 --- Nine Lemma}
\hypertarget{o014.mastery.diagram-chasing.001.ex04}{ }
\label{mastery:diagram-chasing-001-ex04}

Consider the commutative diagram of $R$-modules
\[
\begin{tikzcd}[column sep=small]
	& 0 \arrow[d] & 0 \arrow[d] & 0 \arrow[d] & \\
	0 \arrow[r] & A' \arrow[r, "f'"] \arrow[d, "\alpha"'] & B' \arrow[r, "g'"] \arrow[d, "\beta"'] & C' \arrow[r] \arrow[d, "\gamma"'] & 0 \\
	0 \arrow[r] & A \arrow[r, "f"] \arrow[d, "\overline\alpha"'] & B \arrow[r, "g"] \arrow[d, "\overline\beta"'] & C \arrow[r] \arrow[d, "\overline\gamma"'] & 0 \\
	0 \arrow[r] & A'' \arrow[r, "f''"] \arrow[d] & B'' \arrow[r, "g''"] \arrow[d] & C'' \arrow[r] \arrow[d] & 0 \\
	& 0 & 0 & 0 &
\end{tikzcd}
\]
with all three columns exact. Suppose that the top and middle rows are
exact. Prove that the bottom row is also exact.

% stable-solution-id: o014.mastery.diagram-chasing.001.sol04
\paragraph{Solution.}
\hypertarget{o014.mastery.diagram-chasing.001.sol04}{ }
First we prove that $f''$ is injective. Take $a''\in A''$ with
$f''(a'')=0$, and lift $a''$ to $a\in A$. Since
$\overline\beta f(a)=f''\overline\alpha(a)=0$, there is a $b'\in B'$ with
$\beta(b')=f(a)$. Applying $g$ gives
\[
	\gamma(g'(b'))=g(\beta(b'))=g(f(a))=0.
\]
The morphism $\gamma$ is injective, so $g'(b')=0$. Exactness of the top row
gives $b'=f'(a')$ for some $a'\in A'$. Hence
$f(a-\alpha(a'))=0$. Since $f$ is injective, $a=\alpha(a')$, and thus
$a''=\overline\alpha(a)=0$.

Next, $g''$ is surjective. For $c''\in C''$, choose $c\in C$ mapping to
$c''$. Surjectivity of $g$ gives $b\in B$ with $g(b)=c$; hence
$g''(\overline\beta(b))=\overline\gamma(g(b))=c''$.

It remains to check exactness at $B''$. Commutativity immediately gives
$g''f''=0$. Conversely, take $b''\in B''$ with $g''(b'')=0$ and lift
$b''$ to $b\in B$. Since
$\overline\gamma(g(b))=g''(\overline\beta(b))=0$, there is a $c'\in C'$
with $\gamma(c')=g(b)$. Choose $b'\in B'$ with $g'(b')=c'$. Then
\[
	g(b-\beta(b'))=g(b)-\gamma(g'(b'))=0.
\]
Exactness of the middle row gives $a\in A$ with
$f(a)=b-\beta(b')$. Mapping this equation to the bottom row yields
\[
	f''(\overline\alpha(a))
	=\overline\beta(f(a))
	=\overline\beta(b)=b''.
\]
Thus $\Ker(g'')=\Image(f'')$, and the bottom row is exact.

% stable-exercise-id: o014.mastery.diagram-chasing.001.ex05
\subsubsection*{Exercise 5 --- Constructing the connecting morphism}
\hypertarget{o014.mastery.diagram-chasing.001.ex05}{ }
\label{mastery:diagram-chasing-001-ex05}

Return to the diagram of short exact sequences
\[
\begin{tikzcd}
	0 \arrow[r] & A' \arrow[r, "i"] \arrow[d, "\alpha'"'] & A \arrow[r, "p"] \arrow[d, "\alpha"'] & A'' \arrow[r] \arrow[d, "\alpha''"] & 0 \\
	0 \arrow[r] & B' \arrow[r, "j"'] & B \arrow[r, "q"'] & B'' \arrow[r] & 0.
\end{tikzcd}
\]
For $x''\in\Ker(\alpha'')$, choose $x\in A$ with $p(x)=x''$. Show that
$\alpha(x)=j(y')$ for some $y'\in B'$, and then define
\[
	\delta(x''):=[y']\in\Coker(\alpha')=B'/\Image(\alpha').
\]
Prove that $\delta$ is well defined, is a homomorphism, and is natural with
respect to morphisms between two such diagrams.

% stable-solution-id: o014.mastery.diagram-chasing.001.sol05
\paragraph{Solution.}
\hypertarget{o014.mastery.diagram-chasing.001.sol05}{ }
Since $x''\in\Ker(\alpha'')$,
\[
	q(\alpha(x))=\alpha''(p(x))=\alpha''(x'')=0.
\]
Exactness of the bottom row gives $\alpha(x)\in\Image(j)$, so there is a
$y'\in B'$ with $j(y')=\alpha(x)$. The element $y'$ is unique because $j$
is injective.

If $x_1$ is another lift of $x''$, then
$x_1-x\in\Ker(p)=\Image(i)$; write $x_1-x=i(a')$. If
$j(y_1')=\alpha(x_1)$, commutativity gives
\[
	j(y_1'-y'-\alpha'(a'))
	=\alpha(x_1)-\alpha(x)-\alpha i(a')=0.
\]
Injectivity of $j$ gives $y_1'-y'=\alpha'(a')$. Thus
$[y_1']=[y']$ in $\Coker(\alpha')$, and $\delta$ is independent of the
choice of $x$. For $x_1''$ and $x_2''$, the sum of chosen lifts and the sum
of the corresponding elements $y'$ are valid data for $x_1''+x_2''$;
therefore $\delta$ is additive and $R$-linear.

For naturality, take a morphism from the diagram above to a similar diagram
whose every square commutes. If $x$ and $y'$ are chosen for $x''$, their
images in the second diagram are valid choices for the image of $x''$.
Hence the two paths in the square
\[
\begin{tikzcd}
	\Ker(\alpha'') \arrow[r, "\delta"] \arrow[d] & \Coker(\alpha') \arrow[d] \\
	\Ker(\widetilde\alpha'') \arrow[r, "\widetilde\delta"'] & \Coker(\widetilde\alpha')
\end{tikzcd}
\]
send $x''$ to the same class of the image of $y'$. The square commutes, so
$\delta$ is natural. This is the elementwise version of the canonical
property in Remark~\ref{rem:connecting-canonical}.

% stable-exercise-id: o014.mastery.diagram-chasing.001.ex06
\subsubsection*{Exercise 6 --- Proving the Snake Lemma}
\hypertarget{o014.mastery.diagram-chasing.001.ex06}{ }
\label{mastery:diagram-chasing-001-ex06}

Using the diagram and $\delta$ from Exercise 5, prove that the sequence
\[
\begin{aligned}
	0 \longrightarrow \Ker(\alpha') \longrightarrow \Ker(\alpha)
	&\longrightarrow \Ker(\alpha'') \xrightarrow{\delta}
	\Coker(\alpha') \\
	&\longrightarrow \Coker(\alpha)
	\longrightarrow \Coker(\alpha'') \longrightarrow 0
\end{aligned}
\]
is exact. State clearly where exactness of the two rows is used in the chase.

% stable-solution-id: o014.mastery.diagram-chasing.001.sol06
\paragraph{Solution.}
\hypertarget{o014.mastery.diagram-chasing.001.sol06}{ }
All arrows other than $\delta$ are induced by $i,p,j,q$. The map
$\Ker(\alpha')\to\Ker(\alpha)$ is injective because $i$ is injective.

At $\Ker(\alpha)$, take $x\in\Ker(\alpha)$ that maps to zero in
$\Ker(\alpha'')$. Then $p(x)=0$, so exactness of the top row gives
$x=i(x')$. The equation
$0=\alpha(x)=j\alpha'(x')$ and injectivity of $j$ give
$x'\in\Ker(\alpha')$. The reverse inclusion follows immediately from
$pi=0$. Thus the sequence is exact at $\Ker(\alpha)$.

At $\Ker(\alpha'')$, an element coming from $x\in\Ker(\alpha)$ has
$y'=0$ in the construction of Exercise 5, so it maps to zero under
$\delta$. Conversely, suppose that $x''\in\Ker(\alpha'')$ and
$\delta(x'')=0$. Choose $x\in A$ and $y'\in B'$ as in the construction of
$\delta$. Since $[y']=0$ in $\Coker(\alpha')$, there is an $a'\in A'$ with
$y'=\alpha'(a')$. Therefore
\[
	\alpha(x-i(a'))=j(y')-j\alpha'(a')=0,
	\qquad p(x-i(a'))=x''.
\]
Thus $x''$ comes from $\Ker(\alpha)$.

At $\Coker(\alpha')$, the next arrow sends $[y']$ to $[j(y')]$. If
$[y']=\delta(x'')$ and $x$ is the lift used, then
$j(y')=\alpha(x)$, so $[j(y')]=0$ in $\Coker(\alpha)$. Conversely, if
$[j(y')]=0$, there is an $x\in A$ with $j(y')=\alpha(x)$. Commutativity
gives
\[
	\alpha''p(x)=q\alpha(x)=qj(y')=0.
\]
Thus $p(x)\in\Ker(\alpha'')$, and the construction immediately gives
$\delta(p(x))=[y']$.

At $\Coker(\alpha)$, the image of every $[y']\in\Coker(\alpha')$ maps to
$[qj(y')]=0$. Conversely, suppose that $[b]\in\Coker(\alpha)$ maps to
zero. This means that $q(b)=\alpha''(x'')$ for some $x''\in A''$.
Surjectivity of $p$ gives $x\in A$ with $p(x)=x''$. Then
\[
	q(b-\alpha(x))=q(b)-\alpha''p(x)=0.
\]
Exactness of the bottom row gives $b-\alpha(x)=j(y')$ for some $y'\in B'$.
Since $\alpha(x)$ vanishes in $\Coker(\alpha)$, the class $[b]$ is the
image of $[y']$.

Finally, $\Coker(\alpha)\to\Coker(\alpha'')$ is surjective because $q$
itself is surjective: lift any representative in $B''$ to $B$. Exactness
has now been checked at every term. This is the Snake Lemma
\ref{prop:snake-lemma} for two short exact sequences of modules.

% stable-exercise-id: o014.mastery.diagram-chasing.001.ex07
\subsubsection*{Exercise 7 --- Computing the snake explicitly}
\hypertarget{o014.mastery.diagram-chasing.001.ex07}{ }
\label{mastery:diagram-chasing-001-ex07}

Apply the Snake Lemma to the diagram of abelian groups
\[
\begin{tikzcd}
	0 \arrow[r] & \Z \arrow[r, "\times 6"] \arrow[d, "\times 4"'] & \Z \arrow[r, "\pi"] \arrow[d, "\times 4"'] & \Z/6\Z \arrow[r] \arrow[d, "\times 4"] & 0 \\
	0 \arrow[r] & \Z \arrow[r, "\times 6"'] & \Z \arrow[r, "\pi"'] & \Z/6\Z \arrow[r] & 0.
\end{tikzcd}
\]
Compute all kernels and cokernels, compute $\delta$ on a generator, and then
write down and verify the resulting exact sequence.

% stable-solution-id: o014.mastery.diagram-chasing.001.sol07
\paragraph{Solution.}
\hypertarget{o014.mastery.diagram-chasing.001.sol07}{ }
The first two vertical maps are multiplication by $4$ on $\Z$, so their
kernels are zero and their cokernels are $\Z/4\Z$. In the third column,
\[
	\Ker(\times4:\Z/6\Z\to\Z/6\Z)=\{[0],[3]\}\simeq\Z/2\Z.
\]
The image of multiplication by $4$ on $\Z/6\Z$ is
$\{[0],[2],[4]\}$, so its cokernel is isomorphic to $\Z/2\Z$.

To compute $\delta([3])$, lift $[3]$ to $3\in\Z$. The middle vertical map
sends it to $12$. Since $12=6\cdot2$, the construction of the connecting
morphism gives
\[
	\delta([3])=[2]\in\Z/4\Z.
\]
Under the identification of generators
$[3]\leftrightarrow[1]\in\Z/2\Z$, the map $\delta$ is
$[1]\mapsto[2]$.

The map from the first cokernel to the second is induced by $\times6$, and
hence on $\Z/4\Z$ it equals $\times2$. The next map is reduction modulo $2$.
The Snake Lemma sequence is therefore
\[
	0\longrightarrow0\longrightarrow0\longrightarrow\Z/2\Z
	\mathop{\longrightarrow}^{[1]\mapsto[2]}\Z/4\Z
	\xrightarrow{\times2}\Z/4\Z
	\xrightarrow{\bmod 2}\Z/2\Z\longrightarrow0.
\]
The image of $\delta$ is $\{[0],[2]\}=\Ker(\times2)$, while the image of
$\times2$ is also $\{[0],[2]\}=\Ker(\bmod 2)$. The last map is
surjective, and $\delta$ is injective. Thus the sequence is indeed exact.

% stable-exercise-id: o014.mastery.diagram-chasing.001.ex08
\subsubsection*{Exercise 8 --- The connecting morphism in cohomology}
\hypertarget{o014.mastery.diagram-chasing.001.ex08}{ }
\label{mastery:diagram-chasing-001-ex08}

Let
\[
	0\longrightarrow A^\bullet\xrightarrow{i}B^\bullet
	\xrightarrow{q}C^\bullet\longrightarrow0
\]
be a short exact sequence of cochain complexes of $R$-modules, exact in
every degree. For a class $[c]\in\Hm^n(C^\bullet)$, choose a cocycle
representative $c\in C^n$ and a lift $b\in B^n$ with $q(b)=c$.
\begin{enumerate}[(i)]
	\item Prove that there is an $a\in A^{n+1}$ with
	$i(a)=d_B(b)$, and that $a$ is a cocycle.
	\item Prove that the formula $\partial^n[c]=[a]$ gives a well-defined
	homomorphism
	$\partial^n:\Hm^n(C^\bullet)\to\Hm^{n+1}(A^\bullet)$.
	\item Prove exactness of the segment
	\[
		\Hm^n(A^\bullet)\longrightarrow\Hm^n(B^\bullet)
		\longrightarrow\Hm^n(C^\bullet)\xrightarrow{\partial^n}
		\Hm^{n+1}(A^\bullet)\longrightarrow\Hm^{n+1}(B^\bullet).
	\]
	\item Prove that $\partial^n$ is natural with respect to morphisms of
	short exact sequences of complexes.
\end{enumerate}

% stable-solution-id: o014.mastery.diagram-chasing.001.sol08
\paragraph{Solution.}
\hypertarget{o014.mastery.diagram-chasing.001.sol08}{ }
Since $c$ is a cocycle,
\[
	q(d_Bb)=d_C(qb)=d_Cc=0.
\]
Exactness at $B^{n+1}$ gives
$d_Bb\in\Image(i^{n+1})$, so there is an $a\in A^{n+1}$ with
$i(a)=d_Bb$. Next,
\[
	i(d_Aa)=d_B(i(a))=d_B^2b=0.
\]
Since $i$ is injective in degree $n+2$, we obtain $d_Aa=0$, so $a$ is a
cocycle.

If the lift $b$ is replaced by $b+i(u)$, the corresponding element changes
from $a$ to $a+d_Au$, so the class $[a]$ does not change. If the
representative $c$ is replaced by $c+d_Cv$, choose a lift $w\in B^{n-1}$
of $v$ and use $b+d_Bw$ as a lift of the new representative. Since
$d_B(b+d_Bw)=d_Bb$, the resulting class remains $[a]$. In particular, if
$c$ is a coboundary, we may choose a lift of the form $d_Bw$ and obtain
$a=0$. Thus $\partial^n$ depends only on the cohomology class. Taking sums
of lifts shows that $\partial^n$ is a homomorphism.

We check exactness at each term. At $\Hm^n(B^\bullet)$, the image of a
class coming from $A$ clearly vanishes after applying $q$. Conversely, if
$b\in B^n$ is a cocycle and $[q(b)]=0$, write
$q(b)=d_Cv$ and lift $v$ to $w\in B^{n-1}$. Then
$b-d_Bw\in\Ker(q)=\Image(i)$; write $b-d_Bw=i(a)$. Since $b-d_Bw$ is a
cocycle and $i$ is injective, $a$ is also a cocycle. Thus $[b]$ comes from
$\Hm^n(A^\bullet)$.

At $\Hm^n(C^\bullet)$, the image of a cocycle $b\in B^n$ has $d_Bb=0$,
so its class maps to zero under $\partial^n$. Conversely, if
$\partial^n[c]=0$, the construction gives $i(a)=d_Bb$ with $a=d_Au$ for
some $u\in A^n$. Hence
\[
	d_B(b-i(u))=i(a)-i(d_Au)=0,
	\qquad q(b-i(u))=c.
\]
Thus $[c]$ comes from $\Hm^n(B^\bullet)$.

At $\Hm^{n+1}(A^\bullet)$, the class $\partial^n[c]=[a]$ maps to
$[i(a)]=[d_Bb]=0$. Conversely, suppose that $a\in A^{n+1}$ is a cocycle
and $[i(a)]=0$ in $\Hm^{n+1}(B^\bullet)$. Choose $b\in B^n$ with
$i(a)=d_Bb$. The element $c=q(b)$ is a cocycle because
$d_Cc=q(d_Bb)=q(i(a))=0$, and the construction gives
$\partial^n[c]=[a]$. The requested segment is therefore exact.

Finally, take a morphism between two short exact sequences of complexes. If
$c,b,a$ are the chase data above, their images in the second sequence still
satisfy
\[
	q(b)=c,\qquad i(a)=d_Bb.
\]
Thus applying the morphism first and then forming the connecting morphism
gives the same class as forming $\partial^n$ first and then applying the
induced map on cohomology. Therefore $\partial^n$ is natural.

Carrying out this construction for every $n$ joins the exact segments into
the long exact sequence
\[
	\cdots\to\Hm^n(A^\bullet)\to\Hm^n(B^\bullet)\to\Hm^n(C^\bullet)
	\xrightarrow{\partial^n}\Hm^{n+1}(A^\bullet)\to\Hm^{n+1}(B^\bullet)
	\to\cdots.
\]
The morphism $\partial^n$ is conceptually the same connecting morphism as in
the Snake Lemma: lift a cocycle, map its differential forward, and then
descend the result, which now lies in a kernel, to the left-hand term.
